\documentclass[12pt,a4paper]{report}

\usepackage[utf8]{inputenc}
\usepackage[T1]{fontenc}
\usepackage[spanish]{babel}
\usepackage{listings}
\usepackage[dvipsnames]{xcolor}
\usepackage[lmargin=2.5cm, rmargin=5cm]{geometry}
\usepackage{hyperref}
\usepackage{parskip}
\usepackage{titlesec}
\usepackage{booktabs}
\usepackage{tcolorbox}
\usepackage{amsmath}
\usepackage{array}
\usepackage{lscape}

\definecolor{lisp-bg}{RGB}{248,248,252}
\definecolor{lisp-kw}{RGB}{0,0,180}
\definecolor{lisp-str}{RGB}{160,0,0}
\definecolor{box-blue}{RGB}{240,248,255}
\definecolor{box-blue-border}{RGB}{100,149,237}
\definecolor{box-green}{RGB}{240,255,240}
\definecolor{box-green-border}{RGB}{60,150,60}
\definecolor{box-yellow}{RGB}{255,253,230}
\definecolor{box-yellow-border}{RGB}{180,150,0}

\lstdefinelanguage{DSL}{
  keywords={def-table, libro, hoja, tabla, load, defparameter,
            conditional-rendering, exists, nav, ultima-fila, primera-fila,
            xl-generate, xl-run-generated},
  morekeywords=[2]{countif, counta, sum-range, str, col, trange, param,
                   add, divide, promedio, gt, lt, eq, and, or, not, _if},
  keywordstyle=\color{lisp-kw}\bfseries,
  keywordstyle=[2]\color{box-green-border},
  stringstyle=\color{lisp-str},
  sensitive=false,
  basicstyle=\ttfamily\small,
  backgroundcolor=\color{lisp-bg},
  frame=single,
  framerule=0.4pt,
  rulecolor=\color{gray!40},
  breaklines=true,
  numbers=left,
  numberstyle=\tiny\color{gray},
  numbersep=8pt,
  xleftmargin=20pt,
  aboveskip=6pt,
  belowskip=6pt,
  showstringspaces=false,
  tabsize=2,
}

\lstset{language=DSL}

\tcbuselibrary{skins,breakable}

\newtcolorbox{infobox}[1]{
  colback=box-blue, colframe=box-blue-border,
  fonttitle=\bfseries, title=#1,
  breakable, sharp corners, left=6pt, right=6pt, top=4pt, bottom=4pt
}

\newtcolorbox{advertencia}[1]{
  colback=box-yellow, colframe=box-yellow-border,
  fonttitle=\bfseries, title=#1,
  breakable, sharp corners, left=6pt, right=6pt, top=4pt, bottom=4pt
}

\titleformat{\chapter}[hang]{\large\bfseries}{\thechapter.}{0.5em}{}[\vspace{2pt}\hrule]
\titleformat{\section}[hang]{\normalsize\bfseries}{\thesection}{0.5em}{}
\titleformat{\subsection}[hang]{\normalsize\itshape\bfseries}{\thesubsection}{0.5em}{}

\hypersetup{hidelinks}
\setlength{\parskip}{6pt}

\begin{document}

\chapter{Tutorial básico: tablas, expresiones y resaltado}
\label{cap:tutorial-basico}

Este capítulo presenta los elementos fundamentales del DSL mediante un
ejemplo mínimo que iremos construyendo en tres pasos:

\begin{enumerate}
  \item Crear una tabla vacía con solo los encabezados y exportarla
    (al backend Excel, por ejemplo) para verificar que el sistema
    funciona.
  \item Añadir los datos de entrada.
  \item Incorporar una columna calculada y resaltado condicional.
\end{enumerate}

Al terminar sabrá cómo se definen las columnas de una tabla, cómo se
añaden expresiones que se calculan solas, cómo se resaltan celdas según
una condición y cómo se genera el archivo de salida.

\section{Qué vamos a construir}

Queremos una tabla con cuatro columnas:

\begin{center}
\begin{tabular}{|l|c|c|r|}
\toprule
\textbf{Nombre} & \textbf{Nota 1} & \textbf{Nota 2} & \textbf{Promedio} \\
\midrule
\textcolor{red}{Ana}    & 85 & 90 & 87.5 \\
Luis   & 70 & 75 & 72.5 \\
\textcolor{red}{Elena}  & 92 & 88 & 90.0 \\
Carlos & 60 & 65 & 62.5 \\
\textcolor{red}{Sofía}  & 78 & 82 & 80.0 \\
\textcolor{red}{Pedro}  & 95 & 91 & 93.0 \\
\bottomrule
\end{tabular}
\end{center}

Además:
\begin{itemize}
  \item La columna \textbf{Promedio} no se escribe a mano: se calcula
    automáticamente como el promedio de las dos notas.
  \item Hay un \textbf{número fijo} (80, por ejemplo) almacenado en el
    libro.
  \item Los nombres de las personas cuyo promedio supera ese número
    aparecen \textbf{resaltados en rojo}.
\end{itemize}

\section{Conceptos básicos}

Antes de escribir código, conviene entender la arquitectura del DSL:

\begin{itemize}
  \item Una \textbf{tabla} es una estructura con columnas.  Cada columna
    tiene un nombre interno (símbolo del DSL que se usa en las expresiones del lenguaje)
    y una etiqueta (texto visible del encabezado).  La tabla se define
    con \texttt{def-table}.
  \item Una \textbf{hoja} es una página del libro que agrupa una o más
    tablas.  Se crea con \texttt{hoja}.
  \item Un \textbf{libro} es el contenedor principal que reúne una o más
    hojas.  Se define con \texttt{libro}.
  \item El libro se \textbf{exporta} a un backend concreto (por ejemplo,
    el backend Excel que acompaña al DSL) mediante \texttt{xl-generate}
    y \texttt{xl-run-generated}.
\end{itemize}

Todo el código se escribe en Common Lisp, en un archivo con extensión
\texttt{.lisp}.  El DSL se encarga de traducir esa descripción a un
programa Python que, a través del backend Excel, genera el archivo de
salida.

Para seguir este tutorial cree el archivo \texttt{tutorial-basico.lisp}
dentro de la carpeta \texttt{tutorial-basico/} (donde ya están
\texttt{dsl-directo.lisp} y el resto de los archivos del proyecto).

\medskip
\noindent\textbf{Nota:} a lo largo del tutorial iremos completando el
mismo archivo \texttt{tutorial-basico.lisp}.  Al final tendrá el
programa completo.

%% ===================================================================
\section{Paso 1: Tabla con solo los encabezados}
%% ===================================================================

Vamos a empezar por lo más básico: definir la estructura de la tabla
y crear un libro que la contenga, sin datos todavía.  El objetivo es
confirmar que el pipeline completo funciona (definir, generar,
exportar) y ver cómo se ven los encabezados en el archivo de salida.

\subsection{Definir la tabla}

La macro \texttt{def-table} define una tabla.  Le damos
un nombre y la lista de columnas.  Cada columna se escribe como un par
\texttt{(identificador \textquotedbl Etiqueta\textquotedbl)}, donde \texttt{identificador} es un
símbolo de Common Lisp con el que el DSL reconoce la columna (se usa
en expresiones como \texttt{(col identificador)} para referirse a ella)
y \texttt{\textquotedbl Etiqueta\textquotedbl} es el texto que se
mostrará como encabezado de columna en el archivo de salida.
Ambos pueden ser distintos: el identificador es un símbolo interno del
lenguaje, la etiqueta es externa (lo que ve quien abre el archivo).

\begin{lstlisting}[caption={def-table con etiquetas visibles.}]
(def-table tabla-notas
  ((nombre "Nombre") (nota1 "Nota 1") (nota2 "Nota 2") (promedio "Promedio"))
  :cell-height 1
  :cell-width  1)
\end{lstlisting}

\noindent Aquí:
\begin{itemize}
  \item \texttt{nombre} es el identificador interno de la primera
    columna y \texttt{"Nombre"} es la etiqueta que aparecerá como
    encabezado de columna en el archivo de salida.
  \item \texttt{nota1} y \texttt{nota2} son las columnas de las notas.
  \item \texttt{promedio} es la columna donde más adelante pondremos
    la expresión.
\end{itemize}

La tabla anterior produce un encabezado como este en el archivo de
salida:

\begin{center}
\begin{tabular}{|l|c|c|c|}
\toprule
\textbf{Nombre} & \textbf{Nota 1} & \textbf{Nota 2} & \textbf{Promedio} \\
\midrule
\bottomrule
\end{tabular}
\end{center}

(Por ahora no hay filas de datos; las añadiremos en el paso~2.)

\begin{infobox}{Identificador vs.\ etiqueta}
  El DSL maneja dos nombres por columna:
  \begin{itemize}
    \item El \textbf{identificador} (p.\ ej. \texttt{nombre}) es un
      símbolo de Common Lisp que usa el lenguaje en las expresiones
      como \texttt{(col nombre)}.  Es el que escribe quien programa
      el DSL.
    \item La \textbf{etiqueta} (p.\ ej. \texttt{"Nombre completo"}) es
      el texto que aparece en el encabezado de la columna en el archivo
      de salida.  Es lo que ve quien abre el archivo.
  \end{itemize}
  Pueden ser distintos.  En el ejemplo, \texttt{nombre} es el
  identificador interno y \texttt{"Nombre"} (o \texttt{"Nombre
    completo"}) es la etiqueta visible.
\end{infobox}

\begin{advertencia}{No ejecute este fragmento solo}
  La macro \texttt{def-table} solo tiene sentido dentro del DSL completo.
  Si intenta cargar este fragmento en SBCL sin haber cargado antes
  \texttt{dsl-directo.lisp} obtendrá un error.  No se preocupe: el
  programa completo que verá al final incluye la instrucción
  \texttt{(load ...)} necesaria.
\end{advertencia}

\subsection{Crear el libro y exportar}

Con la tabla definida, creamos un libro con una hoja que contenga la
tabla.  Como aún no tenemos datos, el \texttt{tabla} se escribe sin
\texttt{:data}.

\begin{lstlisting}[caption={Paso 1: libro con tabla vacía.}]
(def-table tabla-notas
  ((nombre "Nombre") (nota1 "Nota 1") (nota2 "Nota 2") (promedio "Promedio"))
  :cell-height 1
  :cell-width  1)

(libro tutorial-basico
  :hojas (list
    (hoja "Notas"
      (tabla tabla-notas)))

(xl-generate tutorial-basico "tutorial-basico.py")
(xl-run-generated "tutorial-basico.py")
\end{lstlisting}

\noindent Ejecute desde la terminal:

\begin{lstlisting}[language=bash, numbers=none, backgroundcolor=\color{gray!10},
                   basicstyle=\ttfamily\small]
cd tutorial-basico/
sbcl --script tutorial-basico.lisp
\end{lstlisting}

Esto carga el DSL, construye el libro, genera \texttt{tutorial-basico.py}
y lo ejecuta para crear \texttt{Archivo-Excel.xlsx} (a través del
backend Excel).  Al abrir el archivo en un programa de planilla verá
una hoja llamada ``Notas'' con los encabezados de las cuatro columnas
y ninguna fila de datos.

\begin{infobox}{¿Por qué sin datos?}
  En este primer paso lo importante es verificar que todo el pipeline
  funciona: el código Lisp se traduce a Python, el Python se ejecuta y
  se produce un archivo de salida (en este caso .xlsx).  Una vez que
  esto funcione, podemos añadir contenido con tranquilidad.
\end{infobox}

%% ===================================================================
\section{Paso 2: Añadir los datos}
%% ===================================================================

Ahora que el sistema funciona, vamos a meter los datos en la tabla.

\subsection{Los datos en Common Lisp}

Los datos se guardan en una variable con \texttt{defparameter}.  Cada
fila es una lista con cuatro elementos (uno por columna).  Las primeras
tres filas se llaman de \textbf{prefijo}: no contienen datos como tal,
sino que definen cómo se ve la tabla antes de los datos.  En concreto:
las dos primeras filas están en blanco (sirven como separadores
visuales en el archivo de salida) y la tercera contiene los encabezados
de columna.  Las filas siguientes (a partir de la cuarta) son los datos
propiamente dichos.

En la lista de datos, las tres primeras filas son de prefijo y el
resto son datos:

\begin{center}
\begin{tabular}{lccc}
\toprule
\textbf{Nombre} & \textbf{Nota 1} & \textbf{Nota 2} & \textbf{Promedio} \\
\midrule
Ana    & 85 & 90 & -- \\
Luis   & 70 & 75 & -- \\
Elena  & 92 & 88 & -- \\
Carlos & 60 & 65 & -- \\
Sofía  & 78 & 82 & -- \\
Pedro  & 95 & 91 & -- \\
\bottomrule
\end{tabular}
\end{center}

\begin{lstlisting}[caption={Datos de entrada.}]
(defparameter *DATOS*
  '(("" "" "" "")
    ("" "" "" "")
    ("Nombre" "Nota 1" "Nota 2" "Promedio")
    ("Ana"     "85" "90" "")
    ("Luis"    "70" "75" "")
    ("Elena"   "92" "88" "")
    ("Carlos"  "60" "65" "")
    ("Sofia"   "78" "82" "")
    ("Pedro"   "95" "91" "")))
\end{lstlisting}

Observe que la columna \texttt{Promedio} se deja con una cadena vacía
(\texttt{""}) en los datos.  En el código de la lista, cada par de
comillas dobles \texttt{""} representa ese valor vacío: es la manera de
decirle al DSL que ``esta celda no tiene valor fijo, la calculará
automáticamente una expresión del lenguaje''.  Nótese que \texttt{""} en Common Lisp
no es lo mismo que la lista vacía \texttt{()}: \texttt{""} es un texto
de longitud cero, mientras que \texttt{()} es una lista sin elementos.

\begin{advertencia}{Formato de las celdas vacías}
  Las celdas que corresponden a columnas calculadas se marcan con
  \texttt{""} (cadena vacía).  Cada par \texttt{""} en la lista de datos
  es un marcador de posición que el generador sustituye por la expresión
  correspondiente al generar el archivo de salida.
\end{advertencia}

\subsection{Modificar el libro para incluir los datos}

Ahora modificamos el programa del paso 1 para que la tabla use los
datos:

\begin{lstlisting}[caption={Paso 2: libro con datos.}]
(defparameter *DATOS*
  '(("" "" "" "")
    ("" "" "" "")
    ("Nombre" "Nota 1" "Nota 2" "Promedio")
    ("Ana"     "85" "90" "")
    ("Luis"    "70" "75" "")
    ("Elena"   "92" "88" "")
    ("Carlos"  "60" "65" "")
    ("Sofia"   "78" "82" "")
    ("Pedro"   "95" "91" "")))

(def-table tabla-notas
  ((nombre "Nombre") (nota1 "Nota 1") (nota2 "Nota 2") (promedio "Promedio"))
  :cell-height 1
  :cell-width  1)

(libro tutorial-basico
  :hojas (list
    (hoja "Notas"
      (tabla tabla-notas
        :data *DATOS*)))

(xl-generate tutorial-basico "tutorial-basico.py")
(xl-run-generated "tutorial-basico.py")
\end{lstlisting}

Ejecute de nuevo:

\begin{lstlisting}[language=bash, numbers=none, backgroundcolor=\color{gray!10},
                   basicstyle=\ttfamily\small]
cd tutorial-basico/
sbcl --script tutorial-basico.lisp
\end{lstlisting}

Ahora el archivo de salida muestra los encabezados (\texttt{Nombre},
\texttt{Nota 1}, \texttt{Nota 2}, \texttt{Promedio}) y debajo las filas
con los datos de cada persona.  La columna \texttt{Promedio} sigue vacía
porque aún no le hemos dicho cómo calcularla.

%% ===================================================================
\section{Paso 3: Columna calculada y resaltado}
%% ===================================================================

En este paso añadimos la expresión que calcula el promedio automáticamente
y la regla que resalta en rojo los nombres de quienes superen la nota
mínima.

\subsection{Columna calculada}

Queremos que los valores de la columna \texttt{promedio} se calculen
automáticamente a partir de \texttt{nota1} y \texttt{nota2}:
\[
\text{promedio} = \frac{\text{nota1} + \text{nota2}}{2}
\]

En lenguaje natural diríamos: ``para cada fila de la tabla, promedia
los valores de las columnas \texttt{nota1} y \texttt{nota2}''.  El DSL
ofrece la forma \texttt{promedio} que hace exactamente eso: recibe una
lista de columnas, genera la suma y la divide entre la cantidad de
columnas automáticamente.

\begin{lstlisting}[caption={Tabla con columna calculada.}]
(def-table tabla-notas
  ((nombre "Nombre") (nota1 "Nota 1") (nota2 "Nota 2") (promedio "Promedio"))
  :cell-height 1
  :cell-width  1
  :computed
    ((promedio (promedio (col nota1) (col nota2)))))
\end{lstlisting}

\noindent Cada pieza significa:

\begin{itemize}
  \item \texttt{(col nota1)} --- toma el valor de la columna
    \texttt{nota1} en la fila actual (la fila que se está procesando
    en ese momento, distinta para cada alumno).
  \item \texttt{(promedio (col nota1) (col nota2))} --- suma nota1 y
    nota2, y divide entre 2 (porque recibe dos argumentos).
\end{itemize}

\begin{infobox}{¿Cómo sabe \texttt{promedio} entre cuánto dividir?}
  \texttt{promedio} cuenta automáticamente los argumentos que recibe.
  Con \texttt{(promedio (col n1) (col n2))} genera
  \texttt{((n1+n2)/2)}; si recibe tres columnas genera
  \texttt{((n1+n2+n3)/3)}.  No hace falta indicar el divisor.
\end{infobox}

\begin{infobox}{Otras funciones disponibles}
  Además de \texttt{promedio}, el DSL ofrece otras funciones para
  columnas calculadas:
  \begin{itemize}
    \item \texttt{(add a b)} --- suma dos valores.
    \item \texttt{(divide a b)} --- divide a entre b.
    \item \texttt{(concat a b)} --- concatena dos textos.
  \end{itemize}
\end{infobox}

\subsection{Resaltado condicional}

Ahora queremos que los nombres de las personas cuyo promedio supere el
mínimo aparezcan en rojo.  En lenguaje natural: ``si el promedio de
esta fila es mayor que nota-minima, colorea el nombre''.

La opción \texttt{:render} de \texttt{def-table} acepta una forma
\texttt{conditional-rendering} que especifica la condición y las
columnas a colorear.

\begin{lstlisting}[caption={Tabla con columna calculada y resaltado.}]
(def-table tabla-notas
  ((nombre "Nombre") (nota1 "Nota 1") (nota2 "Nota 2") (promedio "Promedio"))
  :cell-height 1
  :cell-width  1
  :computed
    ((promedio (promedio (col nota1) (col nota2))))
  :render
    (conditional-rendering
      :condition (gt (col promedio) (param nota-minima))
      :target-columns (nombre)))
\end{lstlisting}

\noindent Al leer la condición en voz alta:

\begin{itemize}
  \item \texttt{gt} significa ``mayor que'' (\emph{greater than}).
  \item \texttt{(col promedio)} es el promedio de la fila actual.
  \item \texttt{(param nota-minima)} es el valor fijo que definiremos.
  \item \texttt{:target-columns (nombre)} indica qué columna se colorea
    cuando la condición se cumple.  Para colorear varias columnas se
    listan los nombres separados por espacio:
    \texttt{:target-columns (nombre promedio)}.
\end{itemize}

\begin{infobox}{Múltiples columnas o condiciones}
  También puede resaltar más de una columna o usar varias condiciones
  en una sola regla.  Por ejemplo:
  \begin{itemize}
    \item Para resaltar tanto nombres como promedios cuando el promedio
      es mayor que 80: \texttt{:target-columns (nombre promedio)}
    \item Para resaltar cuando el promedio es mayor que 80 \textbf{y}
      la nota 1 es mayor que 70: \texttt{(and (gt (col promedio) 80)
      (gt (col nota1) 70))}
  \end{itemize}
  El DSL permite combinar condiciones lógicas con \texttt{and},
  \texttt{or} y \texttt{not} para crear reglas de resaltado complejas.
\end{infobox}

\begin{infobox}{Más operadores de comparación}
  Además de \texttt{gt} ({\itshape greater than}), el DSL dispone de:
  \begin{itemize}
    \item \texttt{(lt a b)} --- verdadero si a es menor que b.
    \item \texttt{(eq a b)} --- verdadero si a es igual a b.
    \item \texttt{(and a b ...)} --- verdadero si todos lo son.
    \item \texttt{(or a b ...)} --- verdadero si al menos uno lo es.
    \item \texttt{(not a)} --- verdadero si a es falso.
  \end{itemize}
  Con ellos puede construir condiciones compuestas, por ejemplo:
  ``resaltar si el promedio es mayor que 80 {\bf y} la nota 1 es menor
  que 90''.
\end{infobox}

\subsection{El parámetro y la generación}

Hemos usado \texttt{(param nota-minima)} en la condición, pero aún no
hemos dicho de dónde sale ese valor.  Los parámetros son valores fijos
con nombre que se definen al instanciar la tabla, mediante la clave
\texttt{:params}.  Cada parámetro es un par \texttt{(nombre valor)}.

\begin{lstlisting}[caption={Libro completo con parámetros.}]
(libro tutorial-basico
  :hojas (list
    (hoja "Notas"
      (tabla tabla-notas
        :data *DATOS*
        :params ((nota-minima 80))))))

(xl-generate tutorial-basico "tutorial-basico.py")
(xl-run-generated "tutorial-basico.py")
\end{lstlisting}

El nombre del archivo de salida lo define internamente el generador
(por defecto \texttt{Archivo-Excel.xlsx}); no hace falta indicarlo
en el DSL.

\begin{infobox}{El concepto de parámetro}
  Un \texttt{param} es un valor fijo con nombre que se escribe una sola
  vez al instanciar la tabla y se referencia desde cualquier expresión
  de esa tabla con \texttt{(param nombre)}.  Su alcance (\emph{scope})
  es la tabla donde se define: cada \texttt{tabla} recibe sus propios
  parámetros y no comparte los de otras tablas.  Es la manera que tiene
  el DSL de representar constantes que puedan cambiar entre ejecuciones
  sin tocar las expresiones.  En este ejemplo \texttt{(param nota-minima)}
  vale 80.

  En la implementación concreta del generador a Excel, los parámetros se
  traducen a celdas fijas de la hoja (por ejemplo \texttt{\$F\$1}) para
  que las expresiones puedan referenciarlos con referencias absolutas.
  El backend de Python elige automáticamente una celda libre de la hoja
  y escribe el valor en ella, sin mezclarlo con las filas de datos de la
  tabla.
\end{infobox}

A continuación se muestra el contenido íntegro del archivo
\texttt{tutorial-basico.lisp} con los tres pasos integrados:

\section{El programa completo}

\begin{lstlisting}[caption={tutorial-basico.lisp --- programa completo.}]
(load (merge-pathnames "dsl-directo.lisp" *load-truename*))

(defparameter *DATOS*
  '(("" "" "" "")
    ("" "" "" "")
    ("Nombre" "Nota 1" "Nota 2" "Promedio")
    ("Ana"     "85" "90" "")
    ("Luis"    "70" "75" "")
    ("Elena"   "92" "88" "")
    ("Carlos"  "60" "65" "")
    ("Sofia"   "78" "82" "")
    ("Pedro"   "95" "91" "")))

(def-table tabla-notas
  ((nombre "Nombre") (nota1 "Nota 1") (nota2 "Nota 2") (promedio "Promedio"))
  :cell-height 1
  :cell-width  1
  :computed
    ((promedio (promedio (col nota1) (col nota2))))
  :render
    (conditional-rendering
      :condition (gt (col promedio) (param nota-minima))
      :target-columns (nombre)))

(libro tutorial-basico
  :hojas (list
    (hoja "Notas"
      (tabla tabla-notas
        :data *DATOS*
        :params ((nota-minima 80))))))

(xl-generate tutorial-basico "tutorial-basico.py")
(xl-run-generated "tutorial-basico.py")
\end{lstlisting}

\section{Ejecución}

Para ejecutar el tutorial completo:

\begin{lstlisting}[language=bash, numbers=none, backgroundcolor=\color{gray!10},
                   basicstyle=\ttfamily\small]
cd tutorial-basico/
sbcl --script tutorial-basico.lisp
\end{lstlisting}

Esto cargará el DSL, construirá el AST, generará
\texttt{tutorial-basico.py} y lo ejecutará para crear
\texttt{Archivo-Excel.xlsx}.

\section{Resultado esperado}

El archivo \texttt{Archivo-Excel.xlsx} contiene una hoja llamada
``Notas'' con la siguiente estructura:

\begin{center}
\begin{tabular}{ll}
\toprule
\textbf{Rango} & \textbf{Contenido} \\
\midrule
Filas 3--9, columnas A--D & Tabla con 6 personas (fila 3: encabezados) \\
Columna D, filas 4--9     & Promedio calculado con expresión \\
Columna A, filas 4--9     & Nombres en rojo si promedio $> 80$ \\
\bottomrule
\end{tabular}
\end{center}

Con los datos de ejemplo, los nombres de Ana, Elena, Sofía y Pedro
aparecen resaltados en rojo porque su promedio supera 80.  Luis y Carlos
quedan con formato normal.

\section{Resumen de formas utilizadas}

\begin{center}
\begin{tabular}{>{\ttfamily}p{5cm}p{8cm}}
\toprule
\textbf{Forma} & \textbf{Para qué sirve} \\
\midrule
(def-table nombre (cols...) opts...) & Define una tabla \\
:computed ((col expr)) & Columna calculada con una expresión \\
(conditional-rendering :condition :target-columns) & Regla de resaltado condicional \\
(promedio col1 col2 ...) & Promedia las columnas indicadas \\
(gt a b) & Verdadero si a es mayor que b \\
(col nombre) & Valor de la columna en la fila actual \\
(param nombre) & Valor de un parámetro fijo \\
:params ((nombre valor)...) & Define parámetros al instanciar la tabla \\
(libro nombre ...) & Define el libro completo \\
(hoja nombre tablas ...) & Define una hoja \\
(tabla definicion :data :params) & Instancia una tabla (definida con def-table) con datos y parámetros \\
(xl-generate wb archivo) & Genera el script Python \\
(xl-run-generated archivo) & Ejecuta el script y produce el archivo de salida (.xlsx) \\
\bottomrule
\end{tabular}
\end{center}

\end{document}
